\documentclass[12pt]{article}
\usepackage[margin=1in]{geometry}
\usepackage{hyperref}
\usepackage{amsmath}
\usepackage{upgreek}

\newcommand{\bigformula}[1]{{\Large \begin{equation} #1 \end{equation}}}

\title{RotateAI Simulator Derivations}
\author{Maxence Morel Dierckx}
\date{\today}

\begin{document}
\maketitle

\section{Disclaimer}
% General overview of MCU, ISA, CPU core benchmarking and power statistics calculations and why a vast majority of ours are absolutely unrelated to the eventual hardware

The simulator runs x86 binaries and measures x86 instruction counts. All derived metrics use these counts as proxies for ARM Cortex-M33 execution. Important disclaimers:

\begin{enumerate}
	\item x86 instructions are not the same as ARM instructions. Different Instruction Set Architectures (ISAs) produce different instruction counts for identical computation.
	\item x86 vectorises float32s to perform 4-8 floating point operations in a single instruction. The M33 architecture does not implement this ``Single Instruction Multiple Data'' (SIMD) technology, so simulator results will heavily undercount the FLOPS of an actual tag.
	\item All hardware-specific constants (voltage, current draw, DMIPS, maximum frequency) are configurable. Actual values vary with peripheral activity, temperature, and process variation.
	\item Binaries that do not produce output on every sample still incur minimal flag-byte I/O overhead per idle sample in the simulator protocol. This overhead would not be present on actual hardware.
\end{enumerate}

The derived values should be interpreted as \textbf{fermi estimates} of eventual characteristics. They have utility as estimates of feasibility and for relative comparison of different model pipelines. They should not be used as engineering specifications.

\section{Variables}

\begin{tabular}{l l l}
	\textbf{Symbol} & \textbf{Description} & \textbf{Source} \\
	\hline
	$N$ & Average instructions per input sample & \verb|perf stat| / sample count \\
	$f_s$ & Sample rate (Hz) & Config \\
	$f_{\max}$ & Maximum operating frequency (MHz) & Config \\
	$D$ & DMIPS per MHz & Config \\
	$\mu$ & Current per MHz ($\upmu$A/MHz) & Config \\
	$V$ & Supply voltage (V) & Config \\
	$I_s$ & Sleep-mode current ($\upmu$A) & Config \\
\end{tabular}

\section{Minimum Operating Frequency}

The Minimum Operating Frequency is the minimum CPU frequency for which the required instructions per inference can be completed in the given sample period (e.g., 200 ms at 5 Hz).

\begin{itemize}
	\item The sample period is the reciprocal of the sample rate $f_s$. 

	\item The Dhrystone Million Instructions Per Second (DMIPS) per MHz (s$^{-1}$MHz$^{-1}$) determines the number of instructions our processor can execute per unit time. We multiply by $10^6$ to convert to s$^{-1}$Hz$^{-1}$.
\end{itemize}

\bigformula{f_{\min} = \frac{N \cdot f_s}{D \cdot 10^6}}
\\
\\
Note that DMIPS is a benchmark-specific metric. \href{https://homepages.cwi.nl/~steven/dry.c}{Dhrystone} is an integer performance benchmark, which means it differs significantly from floating-point-heavy ML inference. Using $D$ as a proxy for the general throughput of the processor is erroneous, but it is the best option STM publishes.

\section{Energy Per Inference}

Energy is the product of voltage, current, and time.

\begin{itemize}
	\item At a given frequency $f$, the MCU draws $\mu \cdot f$ microamps of current.
	\item One inference takes $N / (f \cdot D \cdot 10^6)$ seconds to complete.
\end{itemize}

{\Large
$$E = V \cdot (\mu \cdot f) \cdot \frac{N}{f \cdot D \cdot 10^6}$$
}
\\
\\
The frequency $f$ cancels. Running faster draws more current for less time, and the product is constant within a voltage range. Converting $\mu$ from $\upmu$A to A and simplifying:

\bigformula{E = \frac{V \cdot \mu \cdot N}{D \cdot 10^{12}}}

\section{Duty Cycle}

The duty cycle is the fraction of time the processor spends running inference versus sleeping. This is not measured from the x86 process---it is estimated from $f_{\min}$ and the target's maximum operating frequency $f_{\max}$. Since $f_{\min}$ is the frequency at which inference exactly fills the sample period ($\delta = 1$), running at any higher frequency proportionally reduces the duty cycle:

\bigformula{\delta = \frac{f_{\min}}{f_{\max}}}
\\
\\
A duty cycle above 1 indicates the inference cannot complete within the sample period at $f_{\max}$.

The report also includes the \textbf{output ratio}: the fraction of input samples on which the pipeline produced output. Where $\delta$ measures overall processor utilisation, the output ratio measures how selectively the pipeline schedules inference.

\section{Power Consumption}

Average power has two components: active inference and sleep. Since $N$ is averaged over all input samples (including idle ones), $E \cdot f_s$ already accounts for pipeline selectivity. During idle time $(1 - \delta)$, the MCU draws sleep current $I_s$.

\bigformula{P = E \cdot f_s + V \cdot I_s \cdot 10^{-3} \cdot (1 - \delta)}

\end{document}